\chapter{Математическая статистика}

\textcolor{gray}{\textit{Тут довольно много пространных рассуждений, на тему того, чем занимается
статистика, и между делом вводятся определения.
Мне лень всё это переписывать, там в основном довольно понятные вещи.}}

\section{Построение оценок}

\begin{definition}
    Оценка $ \hat\theta_n $ является \emph{несмещённой} оценкой параметра $ \theta $, если
    $ \mathbb E \hat \theta_n = \theta $.
\end{definition}

\begin{definition}
    Оценка $ \hat\theta_n $ называется \emph{состоятельной}, если
    $ \hat\theta_n \overarr{P} \theta $.
\end{definition}

\begin{definition}
    Оценка называется \emph{эффективной}, если на ней достигается минимум дисперсии.
\end{definition}

\subsection{Метод моментов}

Пусть неизвестный параметр $ \theta $ может быть выражен через параметры исходного распределения,
\ie $ \theta = f(\alpha_1, \dots, \alpha_r) $, где $ \alpha_i = \mathbb E x^i $.
В качестве оценки по методу моментов берётся статистика
$ \hat\theta = f(\hat\alpha_n^{(1)}, \dots, \hat\alpha_n^{(r)}) $,

\subsection{Метод максимального правдоподобия}

Пусть у исходного распределения имеется некоторая плотность распределения $ p(x, \theta) $, параметр
$ \theta $ которой мы оцениваем.
Построим \emph{функцию правдоподобия}
$$ L(\theta, x_1, \dots, x_n) = \prod_{k = 1}^n p(x_k, \theta) $$
В качестве оценки возьмём то значение $ \hat\theta $, которое максимизирует $ L $.

\section{Проверка статистических гипотез}

Пусть имеется две гипотезы о значении некоторого параметра $ a $:
\begin{itemize}
    \item $ H_0 $: $ a = a_0 $;
    \item $ H_1 $: $ a \ne a_0, \quad a < a_0, \quad a > a_0 $.
\end{itemize}
Если принимается решение $ H_0 $, когда на самом деле верно $ H_1 $ "--- это \emph{ошибка первого
рода}.
Другая ошибочная ситуация называется \emph{ошибкой второго рода}.

Фиксируем вероятность $ \alpha $ ошибки первого рода.
Тогда из всех возможных тестов (статистик) надо выбрать те, которые минимизируют вероятность ошибки
второго рода.

\textcolor{gray}{\textit{На этом моменте записи обрываются\dots}}
